\chapter{Graph Algorithms}\label{ch:graphs}

\begin{center}\itshape
Explore as far as you can before turning back, and remember where you have been:\\
the order in which you finish your visits knows the shape of the graph.\\
\upshape(The white-path theorem, in our words)
\end{center}

In 1882, in the first volume of his \emph{R\'ecr\'eations math\'ematiques}, \'Edouard Lucas devoted a chapter to mazes, and in it he reported a method
due to Tr\'emaux, a former student of the \'Ecole polytechnique turned telegraph engineer, for exploring any labyrinth without a map. In our words: mark the
entrance and the exit of every corridor you take, never take a corridor already marked twice, and when a new corridor brings you to a crossing you have already seen,
turn back the way you came. Tr\'emaux's rule is \emph{depth-first search}: go as deep as you can, and back up only when every way forward is used. Ninety years
later, in 1972, Robert Tarjan showed that this old trick, run on a computer and decorated with two numbers per vertex, answers in linear time a long list
of questions that seemed to need much more: Is there a cycle? In which order can these tasks be done? Which vertices can reach each other? Which
cables, if cut, split the network?

A second story belongs to this chapter. One morning in 1956 Edsger Dijkstra, then a young programmer at the Mathematical Centre in Amsterdam, was
shopping in the city with his fiancée. Tired, they sat down on a caf\'e terrace for a cup of coffee, and he wondered how to compute the shortest route
between two Dutch cities, to show off the centre's new computer. He designed the algorithm there, in his head, in about twenty minutes, and published it only
in 1959, in a note of three pages (this is our paraphrase of his account in a 2010 interview). It is now inside every route planner.

This chapter is about questions asked of a graph that is \emph{given}, and must be answered in time proportional to its size. We prove what breadth-first
and depth-first search compute, find strongly connected components in two different ways, keeping a promise made in Tome~C1 to solve $2$-SAT in linear time,
find the weak points of a network, and then turn to weighted graphs: shortest paths with nonnegative weights (Dijkstra), with arbitrary weights (Bellman and Ford),
between all pairs (Floyd and Warshall, and Johnson's trick of changing the weights without changing the answers), and the cost of minimum spanning trees,
whose correctness was proved in Tome~C1.

\begin{think}
You are in a maze with a piece of chalk. Can you guarantee to visit every crossing and come back to the entrance, walking each corridor at most twice?
What should you write on the walls?
\end{think}

\section{Breadth-first search}\label{sec:bfs}

Throughout, a graph $G=(V,E)$ has $n=|V|$ vertices and $m=|E|$ edges, and is stored by \emph{adjacency lists}: for each vertex $u$, the list $\mathrm{Adj}(u)$ of the
vertices $v$ with $uv\in E$ (or $(u,v)\in E$ for a directed graph). This takes $O(n+m)$ words, and ``linear time'' means $O(n+m)$. The \emph{distance} $\delta(s,v)$ is
the least number of edges of a path from $s$ to $v$, and $\infty$ if there is none.

\emph{Breadth-first search} from $s$ sets $d(s)=0$, puts $s$ in a first-in first-out queue $Q$, and, while $Q$ is not empty, removes the vertex $u$ at the head of $Q$ and,
for each $v\in\mathrm{Adj}(u)$ not yet \emph{discovered} (labelled), sets $d(v)=d(u)+1$ and $\pi(v)=u$ and appends $v$ to $Q$. Each vertex enters $Q$ at most once and each
adjacency list is scanned once: the time is $O(n+m)$. The edges $\pi(v)v$ form the \emph{BFS tree}.

\begin{lemma}[the queue]\label{lem:queue}
At every moment, if $Q$ contains $v_1,\dots,v_r$ from head to tail, then $d(v_1)\le d(v_2)\le\dots\le d(v_r)\le d(v_1)+1$. Consequently the vertices enter (and leave) $Q$ in
nondecreasing order of $d$.
\end{lemma}
\begin{proof}
Induction on the operations. Initially $Q=(s)$. Removing $v_1$ keeps the property, since $d(v_r)\le d(v_1)+1\le d(v_2)+1$. When $v$ is appended while $u$ is processed, $u$ was the
head just before its removal, so the property at that time gives $d(v_r)\le d(u)+1=d(v)$ for the current tail, and $d(u)\le d(v_1)$ for the current head, so that $d(v)=d(u)+1\le d(v_1)+1$.
Each appended vertex has $d$ at least that of the tail, which is the second claim.
\end{proof}

\begin{theorem}\label{thm:bfs}
Breadth-first search discovers exactly the vertices reachable from $s$, and for each of them $d(v)=\delta(s,v)$; the path in the BFS tree from $s$ to $v$ is a shortest path.
\end{theorem}
\begin{proof}
A discovered $v$ has $d(v)$ equal to the length of the tree path $s,\dots,\pi(\pi(v)),\pi(v),v$, so $d(v)\ge\delta(s,v)$, and $v$ is reachable. Suppose some reachable $v$ has $d(v)\ne\delta(s,v)$
(including the case where $v$ is never discovered, $d(v)=\infty$), and take one with $\delta(s,v)$ least. Then $v\ne s$; let $u$ be the vertex before $v$ on a shortest path, so $\delta(s,u)=\delta(s,v)-1$
and, by minimality, $d(u)=\delta(s,u)$; in particular $u$ was discovered, hence later removed from $Q$. At that moment, either $v$ was undiscovered, and then it gets $d(v)=d(u)+1=\delta(s,v)$;
or $v$ had already left $Q$, and then $d(v)\le d(u)$ by the lemma; or $v$ was in $Q$, and then $d(v)\le d(u)+1$ by the lemma. In every case $d(v)\le\delta(s,v)$, a contradiction.
\end{proof}

\begin{corollary}\label{cor:bipartite}
A connected graph is bipartite if and only if no edge joins two vertices with the same BFS label. So bipartiteness is decided, and an odd cycle or a $2$-colouring produced,
in linear time.
\end{corollary}
\begin{proof}
For an edge $uv$, $\delta(s,v)\le\delta(s,u)+1$ and symmetrically, so labels of neighbours differ by at most $1$. If no edge joins equal labels, all edges join consecutive levels and the
parity of $d$ is a proper $2$-colouring. If $d(u)=d(v)=k$ for an edge $uv$, the tree paths from $s$ to $u$ and to $v$ followed by $uv$ form a closed walk of length $2k+1$. An odd closed
walk contains an odd cycle: a shortest odd closed walk is a cycle, since a repeated vertex splits it into two shorter closed walks, one of which is odd. A bipartite graph has no odd cycle.
\end{proof}

\section{Depth-first search}\label{sec:dfs}

\emph{Depth-first search} keeps a global clock $t$, starting at $0$, and colours vertices white (undiscovered), gray (in progress) and black (finished):
\begin{center}
\begin{tabular}{l}
$\mathrm{DFS}(u)$:\\
\quad $t\gets t+1$; $d(u)\gets t$; colour $u$ gray;\\
\quad for each $v\in\mathrm{Adj}(u)$:\\
\qquad if $v$ is white: $\pi(v)\gets u$; $\mathrm{DFS}(v)$;\\
\quad $t\gets t+1$; $f(u)\gets t$; colour $u$ black.
\end{tabular}
\end{center}
The main loop calls $\mathrm{DFS}(u)$ for every vertex $u$ still white, in some order. Each vertex is called once and each list scanned once: linear time. The numbers $d(u)<f(u)$ are
the \emph{discovery} and \emph{finishing} times, all distinct, from $1$ to $2n$. The edges $\pi(v)v$ form the \emph{DFS forest}: $v$ is a child of $u$ when $\mathrm{DFS}(v)$ is called from
$\mathrm{DFS}(u)$, so the descendants of $u$ are exactly the vertices whose call starts while $\mathrm{DFS}(u)$ is running, and the gray vertices at any moment are the
calls in progress, which form a path of the forest from a root down to the current vertex.

\begin{theorem}[parenthesis theorem]\label{thm:parenthesis}
For any two vertices $u\ne v$, the intervals $[d(u),f(u)]$ and $[d(v),f(v)]$ are either disjoint or nested, and $[d(v),f(v)]\subset[d(u),f(u)]$ exactly when $v$ is a descendant of $u$.
\end{theorem}
\begin{proof}
Say $d(u)<d(v)$. If $d(v)<f(u)$, then $\mathrm{DFS}(v)$ starts while $\mathrm{DFS}(u)$ runs, so $v$ is a descendant of $u$, and calls are nested: $\mathrm{DFS}(v)$ returns before $\mathrm{DFS}(u)$, so
$f(v)<f(u)$. Otherwise $f(u)<d(v)$ and the intervals are disjoint. Conversely, if $v$ is a descendant of $u$, its call starts during that of $u$, so $d(u)<d(v)<f(u)$.
\end{proof}

Writing ``$($'' at each discovery and ``$)$'' at each finish produces a well-bracketed word, which is where the name comes from. The next result says that the forest is determined
by what is white when a vertex is discovered.

\begin{theorem}[white-path theorem]\label{thm:whitepath}
A vertex $v$ is a descendant of $u$ in the DFS forest if and only if, at time $d(u)$, there is a path from $u$ to $v$ all of whose vertices are white (except $u$, which just turned gray).
\end{theorem}
\begin{proof}
If $v$ is a descendant of $u$, every vertex $w$ of the tree path from $u$ to $v$ is a descendant of $u$, so $d(w)>d(u)$ and $w$ is white at time $d(u)$. Conversely, let $u=w_0,w_1,\dots,w_k=v$
be a white path at time $d(u)$, and suppose that some $w_i$ is not a descendant of $u$; take the least such $i\ge1$. Then $w_{i-1}$ is $u$ or a descendant of $u$, so $f(w_{i-1})\le f(u)$. Since
$w_i$ was white at time $d(u)$, $d(w_i)>d(u)$. When $\mathrm{DFS}(w_{i-1})$ scans the edge to $w_i$, either $w_i$ is white and becomes a child of $w_{i-1}$, or it was already discovered; in both
cases $d(w_i)<f(w_{i-1})\le f(u)$. So $d(u)<d(w_i)<f(u)$, and by the parenthesis theorem $w_i$ is a descendant of $u$, a contradiction.
\end{proof}

In a directed graph, an edge $(u,v)$ is classified at the moment $\mathrm{DFS}(u)$ scans it: a \emph{tree edge} if $v$ is white; a \emph{back edge} if $v$ is gray (then $v$ is an ancestor of $u$, or $u$ itself);
a \emph{forward edge} if $v$ is black and a descendant of $u$; a \emph{cross edge} if $v$ is black and not a descendant of $u$, in which case $f(v)<d(u)$, by the parenthesis theorem, since $v$ finished
before $u$ did and is not nested in it. Figure~\ref{fig:dfs} shows all four kinds. The one fact we shall use again and again is that, for every edge $(u,v)$ that is not a back edge,
\begin{equation}\label{eq:finish}
f(v)<f(u).
\end{equation}
Indeed, if $v$ is white, it becomes a descendant and finishes first; if $v$ is black, it has already finished.

\begin{figure}[ht]
\centering
\fig{ch04-dfs}
\caption{Depth-first search from $a$, scanning neighbours in alphabetical order, then from $h$. Each vertex carries $d/f$. Tree edges are solid, the back edges $(c,a)$ and $(f,d)$
close cycles, $(a,e)$ is a forward edge, and $(h,e)$, $(h,g)$ are cross edges pointing to vertices that finished before $h$ was discovered.}
\label{fig:dfs}
\end{figure}

\begin{theorem}\label{thm:cycle}
A directed graph has a cycle if and only if a depth-first search produces a back edge. If it has none, listing the vertices by decreasing finishing time gives a topological order,
in linear time.
\end{theorem}
\begin{proof}
A back edge $(u,v)$ together with the tree path from $v$ down to $u$ is a cycle. Conversely, given a cycle, let $v$ be its first vertex to be discovered and $u$ its predecessor on the cycle.
At time $d(v)$ the rest of the cycle is a white path from $v$ to $u$, so $u$ is a descendant of $v$ (Theorem~\ref{thm:whitepath}), and when $u$ scans $(u,v)$, $v$ is still gray: a back edge.
Without back edges, \eqref{eq:finish} says that every edge goes from a larger to a smaller finishing time, which is the definition of a topological order.
\end{proof}

This is the linear-time topological sort announced in Tome~C1 (Proposition~2.2 proved that one exists): it is what \texttt{make} does with its dependencies. Notice that
the proof never needed to know \emph{which} DFS was run; any order of the main loop and of the adjacency lists works.

\section{Strongly connected components}\label{sec:scc}

In a directed graph, write $u\leadsto v$ if there is a path from $u$ to $v$. The relation ``$u\leadsto v$ and $v\leadsto u$'' is an equivalence relation; its classes are the
\emph{strongly connected components}. The \emph{condensation} has one vertex per component and an edge $C\to C'$ when some edge goes from $C$ to $C'\ne C$. It is acyclic, since a cycle of
components would let all their vertices reach each other, merging them into one. Figure~\ref{fig:condensation} shows the components of the graph of Figure~\ref{fig:dfs}.

\begin{figure}[ht]
\centering
\fig{ch04-condensation}
\caption{The strongly connected components of the graph of Figure~\ref{fig:dfs}, and its condensation, which is acyclic.}
\label{fig:condensation}
\end{figure}

For a set $C$ of vertices write $f(C)=\max_{v\in C}f(v)$.

\begin{lemma}\label{lem:sccfinish}
After a depth-first search of $G$, if there is an edge from the component $C$ to a different component $C'$, then $f(C)>f(C')$.
\end{lemma}
\begin{proof}
Suppose first that the first vertex discovered in $C\cup C'$ is some $x\in C$. At time $d(x)$ every vertex of $C\cup C'$ is white and reachable from $x$ along vertices of $C\cup C'$ (inside $C$, then
through the edge into $C'$, then inside $C'$). By the white-path theorem they are all descendants of $x$, so they finish before $x$, and $f(C)=f(x)>f(C')$.

Otherwise the first one is some $y\in C'$. As before, all of $C'$ consists of descendants of $y$, so $f(C')=f(y)$. No vertex of $C$ is reachable from $y$, for otherwise $C'\leadsto C$ and $C\leadsto C'$ would
merge the two components. A vertex of $C$ discovered during $\mathrm{DFS}(y)$ would be a descendant of $y$ by the parenthesis theorem, hence reachable from it; and none was discovered before $d(y)$.
So all of $C$ is discovered after $f(y)$, and $f(C)>f(y)=f(C')$.
\end{proof}

So the vertex finishing last lies in a component with no incoming edge from other components. The \emph{transpose} $G^{\mathsf T}$ has all edges reversed; it has the same components,
but the edges of the condensation are reversed. That suggests a beautiful algorithm.

\begin{theorem}[Kosaraju, 1978; Sharir, 1981]\label{thm:kosaraju}
Run a depth-first search of $G$, recording finishing times. Then run a depth-first search of $G^{\mathsf T}$ whose main loop considers the vertices by decreasing $f$. The trees of the
second search are exactly the strongly connected components, and they appear in a topological order of the condensation of $G$. The total time is linear.
\end{theorem}
\begin{proof}
We show by induction on $k$ that the first $k$ trees of the second search are components. Then the vertices still white form a union of components. The next root $r$ is the white vertex with the largest
$f$; let $C$ be its component, so all of $C$ is white and $f(C)=f(r)$. In $G^{\mathsf T}$, $C$ is still strongly connected, so by the white-path theorem the tree of $r$ contains $C$. An edge of $G^{\mathsf T}$
leaving $C$ goes to a component $C'$ having an edge to $C$ in $G$, so $f(C')>f(C)=f(r)$ by Lemma~\ref{lem:sccfinish}, and since $r$ has the largest $f$ among white vertices, $C'$ is not white: it
was already output. Thus the search from $r$ cannot leave $C$, and the tree is exactly $C$. The same remark shows that every edge of $G$ entering $C$ comes from a component output earlier,
which is the claim about topological order.
\end{proof}

Tarjan had found in 1972 a way to do it with one search only, and this is the algorithm Tome~C1 promised. It pushes each vertex on a stack $S$ when it is discovered, and computes
\[
\mathrm{low}(u)=\min\big(d(u),\ \mathrm{low}(c),\ d(w)\big),
\]
the minimum being over the children $c$ of $u$ and over the edges $(u,w)$ such that $w$ is on $S$ when the edge is scanned.
When $\mathrm{DFS}(u)$ ends, if $\mathrm{low}(u)=d(u)$, it pops $S$ down to $u$ included and outputs the popped vertices as a component. For a component $C$, let its \emph{root} $\rho(C)$ be its first discovered vertex;
by the white-path theorem all of $C$ consists of descendants of $\rho(C)$, so $\rho(C(u))$ is an ancestor of $u$ (or $u$ itself) for every $u$.

\begin{theorem}[Tarjan, 1972]\label{thm:tarjan}
Tarjan's algorithm outputs every strongly connected component exactly once, at the moment its root finishes, and runs in linear time.
\end{theorem}
\begin{proof}
We prove, by induction on time, that up to now each component was output exactly when its root finished, so that (I) a vertex is on $S$ if and only if it has been discovered and the root of
its component has not finished. Granting (I) until $\mathrm{DFS}(u)$ ends, we show $\mathrm{low}(u)=d(u)$ if and only if $u$ is a root, and that then $S$ holds $C(u)$ above $u$.

\emph{(a) $\mathrm{low}(u)\ge d(\rho(C(u)))$.} Apart from $d(u)$, the values entering $\mathrm{low}(u)$ (unfolding the recursion over children) are $d(w)$ for edges $(x,w)$ with $x$ a descendant of $u$ (or $u$)
and $w$ on $S$ when the edge is scanned. By (I) the root $\rho=\rho(C(w))$ is then gray, hence an ancestor of the current vertex $x$. If $\rho$ is an ancestor of $u$, then $\rho\leadsto u\leadsto x\to w\leadsto\rho$,
so $u\in C(w)$ and $d(w)\ge d(\rho)=d(\rho(C(u)))$. Otherwise $\rho$ is a proper descendant of $u$ and $d(w)\ge d(\rho)>d(u)\ge d(\rho(C(u)))$.

\emph{(b) If $u$ is not a root, $\mathrm{low}(u)<d(u)$.} Then $\rho=\rho(C(u))$ is a proper ancestor of $u$. Follow a path from $u$ to $\rho$ and let $(x,w)$ be its first edge leaving the subtree of $u$. All vertices of the path
lie in $C(u)$ (they are on the cycle $u\leadsto\rho\leadsto u$). Since $x$ scans the edge before it finishes and $w$ does not become its child, $w$ was discovered before $f(x)\le f(u)$; not being a descendant of $u$, it
was discovered before $u$ (parenthesis theorem). At that moment $\rho$, an ancestor of $u$, is gray, so $w$ is on $S$ by (I), and $\mathrm{low}(x)\le d(w)<d(u)$. Climbing the tree from $x$ to $u$, $\mathrm{low}(u)\le\mathrm{low}(x)$.

If $u$ is a root, (a) gives $\mathrm{low}(u)\ge d(u)$, so $\mathrm{low}(u)=d(u)$; with (b), the test is exact. Finally let $u$ be a root at time $f(u)$. The vertices above $u$ on $S$ were discovered after $u$, and before $f(u)$
since they were pushed before $u$ finished: they are descendants of $u$. A descendant $v$ of $u$ has its root either in the subtree of $u$, or it would be a proper ancestor of $u$, putting $u$ on a cycle
through $v$ and making that root the root of $C(u)=C(v)$, impossible. Roots that are proper descendants of $u$ have finished and their components were removed; what remains above $u$ is exactly
the descendants whose root is $u$, which is $C(u)$. Popping them keeps (I). Each vertex is pushed and popped once: linear time.
\end{proof}

Now we keep the promise of Tome~C1, Chapter~1. A set of $2$-clauses on variables $x_1,\dots,x_n$ has an \emph{implication graph} on the $2n$ literals, with edges $\neg a\to b$ and $\neg b\to a$ for each clause $a\vee b$.
Tome~C1, Theorem~1.17 (Aspvall, Plass and Tarjan) says that the clauses are satisfiable if and only if no variable has $x$ and $\neg x$ in the same strongly connected component.

\begin{theorem}\label{thm:twosatlinear}
$2$-SAT is solvable in linear time. If no $x$ shares a component with $\neg x$, number the components in a topological order of the condensation, write $t(\ell)$ for the number of the component of the
literal $\ell$, and set $x$ true exactly when $t(x)>t(\neg x)$; this satisfies every clause.
\end{theorem}
\begin{proof}
The graph has $2n$ vertices and twice as many edges as clauses; Theorem~\ref{thm:kosaraju} or \ref{thm:tarjan} gives the components and a topological order in linear time. With this rule, a literal $\ell$ (positive or negative)
is true exactly when $t(\ell)>t(\neg\ell)$. Along any edge $\ell\to\ell'$, $t(\ell)\le t(\ell')$, with equality inside a component. Suppose the clause $a\vee b$ is false: $t(a)<t(\neg a)$ and $t(b)<t(\neg b)$. The edges $\neg a\to b$
and $\neg b\to a$ give
\[
t(a)<t(\neg a)\le t(b)<t(\neg b)\le t(a),
\]
which is absurd.
\end{proof}

Tarjan's algorithm outputs components in \emph{reverse} topological order (a component is output only after everything it can reach), so with it the rule reads: $x$ is true when its component is output first.

\section{Cut vertices and bridges}\label{sec:cutvertices}

Now let $G$ be undirected, connected and simple. A \emph{cut vertex} is a vertex whose removal disconnects $G$; a \emph{bridge} is an edge whose removal disconnects it. In an undirected depth-first
search, every edge is scanned twice, once from each end; an edge is a tree edge if it is the one through which its lower end was discovered, and a \emph{back edge} otherwise.

\begin{lemma}\label{lem:noncross}
In a depth-first search of an undirected graph, the ends of every edge are an ancestor and a descendant.
\end{lemma}
\begin{proof}
Let $uv$ be an edge with $d(u)<d(v)$. When $u$ scans it, before $f(u)$, $v$ is either white, and becomes a child, or already discovered; in both cases $d(u)<d(v)<f(u)$, and the parenthesis theorem applies.
\end{proof}

Write $T_v$ for the subtree of $v$, and let $\mathrm{low}(v)$ be the least of $d(v)$ and the $d(w)$ over back edges $xw$ with $x\in T_v$ (not counting the tree edge from $v$ to its parent). By the lemma, an edge
with one end in $T_v$ and the other outside is either the tree edge from $v$ to its parent or a back edge to a proper ancestor of $v$. The values satisfy
$\mathrm{low}(v)=\min\big(d(v),\min_{c}\mathrm{low}(c),\min_{\text{back edges }vw}d(w)\big)$, over the children $c$ of $v$, and are computed during the search in linear time.

\begin{theorem}[Hopcroft and Tarjan, 1973]\label{thm:cutvertices}
In a depth-first search of a connected simple graph:
\begin{enumerate}
\item the root is a cut vertex if and only if it has at least two children;
\item another vertex $u$ is a cut vertex if and only if it has a child $v$ with $\mathrm{low}(v)\ge d(u)$;
\item a tree edge $uv$, with $u$ the parent, is a bridge if and only if $\mathrm{low}(v)>d(u)$; back edges are never bridges.
\end{enumerate}
So all cut vertices and bridges are found in linear time.
\end{theorem}
\begin{proof}
(1) By Lemma~\ref{lem:noncross} there is no edge between the subtrees of two children of the root, so with two children the root separates them. With one child $c$, $G$ minus the root contains the tree $T_c$, which is connected.

(2) For a back edge $xw$ leaving $T_v$, $w$ is a proper ancestor of $v$, and $d(w)\le d(u)$ exactly when $w$ is $u$ or above $u$. If $\mathrm{low}(v)\ge d(u)$, every edge leaving $T_v$ goes to $u$; removing $u$ separates
$T_v$ from the root. If instead every child $v$ has $\mathrm{low}(v)<d(u)$, each $T_v$ has an edge to a proper ancestor of $u$; and the tree minus $T_u$ is connected and contains the root. So $G-u$ is connected.

(3) A back edge lies on a cycle with the tree path between its ends, so it is not a bridge. If $\mathrm{low}(v)>d(u)$, the only edge leaving $T_v$ is $uv$, a bridge. If $\mathrm{low}(v)\le d(u)$, a back edge goes from $T_v$ to $u$ or
above it, and with tree paths it makes a cycle through $uv$.
\end{proof}

\section{Shortest paths}\label{sec:shortest}

Now edges carry weights $w(u,v)\in\mathbb R$, the weight of a path is the sum of its edge weights, and $\delta(s,v)$ is the infimum of the weights of walks from $s$ to $v$ ($+\infty$ if there is none).
If a cycle of negative weight is reachable from $s$ and reaches $v$, going around it makes $\delta(s,v)=-\infty$. If not, a walk containing a cycle can drop it without getting heavier, so $\delta(s,v)$ is
the minimum over \emph{simple} paths, of which there are finitely many: it is attained, by a path with at most $n-1$ edges. Two remarks will be used constantly: a part of a shortest path is a shortest path between
its ends (otherwise replace it), and $\delta(s,v)\le\delta(s,u)+w(u,v)$ for every edge (the \emph{triangle inequality}).

All algorithms below maintain estimates $D(v)$, starting from $D(s)=0$ and $D(v)=\infty$ otherwise, and change them only by \emph{relaxing} an edge $(u,v)$:
\[
\text{if } D(u)+w(u,v)<D(v),\ \text{set } D(v)\gets D(u)+w(u,v) \text{ and } \pi(v)\gets u.
\]

\begin{lemma}[relaxation]\label{lem:relax}
(i) At all times, each finite $D(v)$ is the weight of some walk from $s$ to $v$, so $D(v)\ge\delta(s,v)$; and $D(v)$ never increases.
(ii) If $s=v_0,v_1,\dots,v_k$ is a path and its edges are relaxed in the order $(v_0,v_1),\dots,(v_{k-1},v_k)$, possibly with other relaxations in between, then afterwards $D(v_k)$ is at most the weight of the path,
and equal to $\delta(s,v_k)$ if the path is a shortest path.
\end{lemma}
\begin{proof}
(i) By induction: a relaxation sets $D(v)$ to the weight of a walk to $u$ extended by $(u,v)$. (ii) By induction on $i$, $D(v_i)$ is at most the weight of the prefix up to $v_i$ once $(v_{i-1},v_i)$ has been relaxed: $D(v_i)\le D(v_{i-1})+w(v_{i-1},v_i)$, and values never increase. For a shortest path
the weight of the prefix is $\delta(s,v_i)$, since a prefix of a shortest path is shortest; with (i), equality.
\end{proof}

\subsection*{Nonnegative weights: Dijkstra}

\emph{Dijkstra's algorithm} keeps the vertices not yet \emph{settled} in a priority queue keyed by $D$; it repeatedly extracts the vertex $u$ of least $D$, declares it settled, and relaxes all edges leaving $u$.

\begin{theorem}[Dijkstra, 1959]\label{thm:dijkstra}
If all weights are nonnegative, every vertex $u$ has $D(u)=\delta(s,u)$ when it is extracted. With a binary heap, the algorithm runs in time $O((n+m)\log n)$.
\end{theorem}
\begin{proof}
Suppose not, and let $u$ be the first vertex extracted with $D(u)\ne\delta(s,u)$; so $D(u)>\delta(s,u)$, and $\delta(s,u)<\infty$. Then $u\ne s$, since $s$ is extracted first with $D(s)=0=\delta(s,s)$. Take a shortest path from $s$ to $u$, and let
$(x,y)$ be its first edge with $x$ settled and $y$ not, just before $u$ is extracted (it exists since $s$ is settled and $u$ is not). When $x$ was extracted, $D(x)=\delta(s,x)$ and $(x,y)$ was relaxed, so $D(y)\le\delta(s,x)+w(x,y)=\delta(s,y)$, the prefix of a shortest path being shortest; with Lemma~\ref{lem:relax}(i), $D(y)=\delta(s,y)$.
Weights being nonnegative, $\delta(s,y)\le\delta(s,u)$. Hence $D(y)\le\delta(s,u)<D(u)$, while $u$ was extracted with the least key among unsettled vertices, $y$ among them. Contradiction.

Each vertex is inserted and extracted once, and each edge relaxed once; a relaxation lowers a key, which is a sift-up from the vertex's position in the heap (the heap stores, for each vertex, its index). Each
operation costs $O(\log n)$ by Proposition~\ref{prop:heapops}.
\end{proof}

\begin{remark}
With the Fibonacci heaps of Fredman and Tarjan (1987), where lowering a key costs $O(1)$ amortised, Dijkstra's algorithm runs in $O(m+n\log n)$. A simpler trick, used in our notebook, skips decrease-key: push
$(D(v),v)$ again after each relaxation and ignore stale entries when extracted. The heap holds at most $m+1$ entries, and $\log m\le2\log n$, so the bound $O((n+m)\log n)$ survives.
\end{remark}

\begin{think}
Where exactly did the proof use $w\ge0$? Draw a graph with three vertices and one negative edge on which Dijkstra's algorithm settles a vertex too early (Exercise~\ref{ex:dijkstraneg}).
\end{think}

\subsection*{Arbitrary weights: Bellman and Ford}

The \emph{Bellman--Ford algorithm} does $n-1$ \emph{rounds}, each relaxing every edge once, in any order; then it makes one more pass, and reports a negative cycle if some edge can still be relaxed.

\begin{theorem}[Bellman, 1958; Ford, 1956; Moore, 1959]\label{thm:bellmanford}
After the $n-1$ rounds, if no negative cycle is reachable from $s$, then $D(v)=\delta(s,v)$ for all $v$, and no edge can be relaxed; if a negative cycle is reachable from $s$, some edge can still be relaxed.
The time is $O(nm)$.
\end{theorem}
\begin{proof}
Without reachable negative cycle, each reachable $v$ has a shortest path with $k\le n-1$ edges, and round $i$ relaxes its $i$-th edge after round $i-1$ relaxed the $(i-1)$-th: by Lemma~\ref{lem:relax}(ii), $D(v)=\delta(s,v)$
after round $k$, and it stays so. Unreachable vertices keep $D=\infty=\delta$. Then the triangle inequality, which holds for $\delta$, says that no relaxation applies (an edge $(u,v)$ with $D(u)=\infty$ cannot be relaxed).

Now let $v_0,v_1,\dots,v_k=v_0$ be a cycle reachable from $s$ with $\sum_iw(v_{i-1},v_i)<0$. After the rounds, all its vertices have finite $D$: each is reached from $s$ by a simple path, of at most $n-1$ edges, whose weight bounds $D$ by Lemma~\ref{lem:relax}(ii). If no edge could be relaxed, $D(v_i)\le D(v_{i-1})+w(v_{i-1},v_i)$ for all $i$, and
summing around the cycle, where both sums of $D$ are the same,
\[
0\le\sum_{i=1}^kw(v_{i-1},v_i)<0.
\]
\end{proof}

Negative weights are not exotic: if a unit of currency $i$ buys $r_{ij}$ units of currency $j$, a sequence of exchanges multiplies the money by the product of the rates, and the weights $-\log r_{ij}$ turn a
profitable cycle, an \emph{arbitrage}, into a negative cycle (Exercise~\ref{ex:arbitrage}).

\subsection*{All pairs: Floyd and Warshall, and Johnson}

Number the vertices $1,\dots,n$ and let $D_k(i,j)$ be the least weight of a path from $i$ to $j$ whose \emph{intermediate} vertices (all but the ends) are in $\{1,\dots,k\}$. So $D_0(i,j)=w(i,j)$ for an edge,
$D_0(i,i)=0$, and $D_0(i,j)=\infty$ otherwise.

\begin{theorem}[Floyd, 1962; Warshall, 1962]\label{thm:floyd}
If $G$ has no negative cycle,
\[
D_k(i,j)=\min\big(D_{k-1}(i,j),\ D_{k-1}(i,k)+D_{k-1}(k,j)\big),
\]
and $D_n(i,j)=\delta(i,j)$. Computing the $n$ matrices in place, with one $n\times n$ array, takes $O(n^3)$ time.
\end{theorem}
\begin{proof}
As before, without negative cycles a walk can shed its cycles, so $D_k(i,j)$ is also the least weight of a \emph{walk} with intermediate vertices in $\{1,\dots,k\}$, attained by a simple path $p$. If $k$ is not on $p$, its
weight is $D_{k-1}(i,j)$. If it is, $k$ occurs once and cuts $p$ into paths $i\to k$ and $k\to j$ with intermediates in $\{1,\dots,k-1\}$, so $w(p)\ge D_{k-1}(i,k)+D_{k-1}(k,j)$. Conversely, both terms of the minimum are weights of
walks with intermediates in $\{1,\dots,k\}$. For $k=n$ all vertices are allowed. In round $k$, row $k$ and column $k$ do not change, because $D_{k-1}(k,k)=0$ (it is at most $0$, and a negative value would be a negative closed walk),
so overwriting entries in place is harmless.
\end{proof}

Exercise~\ref{ex:fwneg} shows that negative cycles are detected on the diagonal. For sparse graphs, $n$ runs of Dijkstra would be faster, if only the weights were nonnegative. Johnson's idea is to change the weights
without changing which paths are shortest. For any function $h:V\to\mathbb R$, put
\[
w_h(u,v)=w(u,v)+h(u)-h(v).
\]
Along a path $p$ from $s$ to $t$ the sum telescopes, $w_h(p)=w(p)+h(s)-h(t)$: all paths between the same ends shift by the same amount, so shortest paths are the same, and cycles keep their weight.

\begin{theorem}[Johnson, 1977]\label{thm:johnson}
Let $G$ have no negative cycle. Add a vertex $q$ with an edge of weight $0$ to every vertex, and let $h(v)=\delta(q,v)$, computed by Bellman--Ford. Then $w_h\ge0$, and running Dijkstra's algorithm from every vertex
with the weights $w_h$ gives $\delta(u,v)=\delta_h(u,v)-h(u)+h(v)$ for all pairs, in time $O(nm+n(n+m)\log n)$.
\end{theorem}
\begin{proof}
No edge enters $q$, so no new cycle appears, and $h$ is finite since every vertex is reachable from $q$. The triangle inequality $\delta(q,v)\le\delta(q,u)+w(u,v)$ is exactly $w_h(u,v)\ge0$. The formula for $\delta$ is the
telescoping identity applied to shortest paths. The time is one Bellman--Ford and $n$ Dijkstras.
\end{proof}

\section{Minimum spanning trees, at a price}\label{sec:mst}

Tome~C1 proved \emph{why} the greedy algorithms for minimum spanning trees are correct: Kruskal's rule is the greedy algorithm of a matroid (Corollary~10.9), and the cut property (Exercise~10.3) justifies Prim's
and Bor\r{u}vka's rules when costs are distinct, which we may assume by breaking ties with the edge numbers. Here we only count. Let $G$ be connected, so $m\ge n-1$.

\begin{proposition}\label{prop:mstcost}
Kruskal's algorithm with union--find, Prim's algorithm with a binary heap, and Bor\r{u}vka's algorithm all run in time $O(m\log n)$.
\end{proposition}
\begin{proof}
\emph{Kruskal}: this is Exercise~\ref{ex:kruskaltime}; $\log m\le2\log n$. \emph{Prim} is Dijkstra's algorithm with the key of an outside vertex $v$ being the cost of the cheapest edge from the tree to $v$, instead of a distance: the
same bookkeeping, $O((n+m)\log n)=O(m\log n)$. \emph{Bor\r{u}vka} works in phases: label the components of the current forest by a search, in $O(n+m)$; scan all edges to find, for each component, its cheapest leaving edge;
add all these edges. They all belong to the minimum spanning tree, by the cut property, so they create no cycle. Each component is joined to at least one other, so the number of components is at least halved by every
phase, and there are at most $\lceil\log_2n\rceil$ phases of linear cost.
\end{proof}

\begin{remark}
Combining Bor\r{u}vka phases with random sampling, Karger, Klein and Tarjan (1995) found a minimum spanning tree in expected linear time. Deterministically, Chazelle (2000) reached $O(m\,\alpha(m,n))$ with the inverse Ackermann
function of Chapter~\ref{ch:datastructures}, and Pettie and Ramachandran (2002) gave an algorithm that is provably optimal, although nobody knows its running time.
\end{remark}

\begin{beauty}
Depth-first search is one recursive procedure of four lines, and yet two integers per vertex, the time you arrive and the time you leave, encode nearly everything about a
directed graph: its cycles are the back edges, a topological order is the reversed order of departures, and the components are the trees of a second search run against the
arrows in that order. Shortest paths have a different beauty: a single local move, relaxation, and a single inequality, $\delta(s,v)\le\delta(s,u)+w(u,v)$, are behind
Dijkstra, Bellman--Ford, Floyd--Warshall, and Johnson's reweighting, where a potential borrowed from physics makes all weights nonnegative without changing a single answer.
\end{beauty}

\section{Exercises}

\begin{exercise}[$\star$]\label{ex:treediam}
In a tree, run BFS from any vertex $s$ and let $a$ be a vertex farthest from $s$; run BFS from $a$ and let $b$ be farthest from $a$. Show that $\delta(a,b)$ is the diameter (the largest distance).
\end{exercise}

\begin{exercise}[$\star\star$]\label{ex:euler}
\emph{Hierholzer's algorithm.} In a connected graph with all degrees even, walk from a vertex along unused edges until stuck; then, while some vertex of the closed walk has an unused edge, start a new walk there and splice it in.
Show that you can only get stuck at the starting vertex, and implement the method in linear time with a stack, yielding an Euler circuit.
\end{exercise}

\begin{exercise}[$\star$]\label{ex:countpaths}
Modify breadth-first search to count, for every $v$, the number of shortest paths from $s$ to $v$, in linear time (counting arithmetic operations as unit cost).
\end{exercise}

\begin{exercise}[$\star$]\label{ex:dijkstraneg}
Give a graph with three vertices and one negative edge on which Dijkstra's algorithm settles a vertex with a wrong distance. Show that adding the same constant to all weights does not repair it.
\end{exercise}

\begin{exercise}[$\star$]\label{ex:arbitrage}
Given exchange rates $r_{ij}>0$ between $n$ currencies, decide in time $O(n^3)$ whether some cycle of exchanges $i_1\to i_2\to\dots\to i_k\to i_1$ multiplies the money by more than $1$.
\end{exercise}

\begin{exercise}[$\star$]\label{ex:dagsp}
Show that in a directed acyclic graph, with arbitrary real weights, shortest paths from $s$ are found in linear time by relaxing the edges leaving each vertex in a topological order. Deduce the longest paths
(the \emph{critical path} of a project whose tasks have durations and precedence constraints).
\end{exercise}

\begin{exercise}[$\star\star$]\label{ex:widest}
The \emph{width} of a path is the least capacity of its edges. Adapt Dijkstra's algorithm to find, from $s$, paths of maximum width to all vertices, and prove it correct. Show also that, in an undirected graph, the path in a
\emph{maximum} spanning tree is a widest path.
\end{exercise}

\begin{exercise}[$\star\star$]\label{ex:robbins}
\emph{Robbins' theorem} (1939). Show that a connected undirected graph can be oriented to become strongly connected if and only if it has no bridge, and find such an orientation in linear time by orienting tree edges
away from the root and back edges towards it.
\end{exercise}

\begin{exercise}[$\star$]\label{ex:fwneg}
Run the recurrence of Theorem~\ref{thm:floyd} on any directed graph. Show that $G$ has a negative cycle if and only if some diagonal entry $D_n(i,i)$ is negative.
\end{exercise}

\begin{exercise}[$\star\star$]\label{ex:zeroone}
If all weights are $0$ or $1$, show that shortest paths from $s$ can be found in linear time, with a double-ended queue in place of the heap: put $v$ at the front after relaxing a $0$-edge and at the back after a $1$-edge.
\end{exercise}

\section{Hints}

\begin{hint}{ex:treediam}
Let $x,y$ be the ends of a longest path. Show that $a$ is as far from $x$ or from $y$ as the other end, using that paths in a tree are unique.
\end{hint}
\begin{hint}{ex:euler}
When the walk enters a vertex $v\ne$ start, it has used an odd number of edges at $v$.
\end{hint}
\begin{hint}{ex:countpaths}
$N(v)$ is the sum of $N(u)$ over the neighbours $u$ with $d(u)=d(v)-1$.
\end{hint}
\begin{hint}{ex:dijkstraneg}
Edges $s\to a$ of weight $2$, $s\to b$ of weight $3$, $b\to a$ of weight $-2$. For the constant, compare paths with different numbers of edges.
\end{hint}
\begin{hint}{ex:arbitrage}
$\prod r>1$ if and only if $\sum(-\log r)<0$. Use Floyd--Warshall and Exercise~\ref{ex:fwneg}, or Bellman--Ford from an added source.
\end{hint}
\begin{hint}{ex:dagsp}
In a topological order, the edges of every path appear in path order. Negate the weights for longest paths.
\end{hint}
\begin{hint}{ex:widest}
Replace $D(u)+w(u,v)$ by $\min(W(u),c(u,v))$ and extract the maximum. For the second part, use the cycle property of the maximum spanning tree.
\end{hint}
\begin{hint}{ex:robbins}
Use Theorem~\ref{thm:cutvertices}(3) to show that every vertex can climb to a proper ancestor.
\end{hint}
\begin{hint}{ex:fwneg}
Prove that $D_k(i,j)\le w(p)$ for every path $p$ from $i$ to $j$ whose intermediate vertices are distinct, in $\{1,\dots,k\}$, and different from $i,j$; apply it to a negative cycle through its largest vertex.
\end{hint}
\begin{hint}{ex:zeroone}
Show that the deque has the property of Lemma~\ref{lem:queue}, then follow the proof of Theorem~\ref{thm:dijkstra}.
\end{hint}

\section{Solutions}

\begin{solution}{ex:treediam}
\idea A vertex farthest from anywhere is an end of some longest path, because a tree has nowhere else to put far vertices.
\writeup Let $P$ be a path from $x$ to $y$ of maximum length $L$. In a tree, the path from any vertex $z$ to $P$ first meets $P$ at some vertex $m_z$, and then the paths from $z$ to $x$ and to $y$ both go through $m_z$ (the walk $z\to m_z$ followed by
the part of $P$ has no repeated vertex, so it is \emph{the} path), so $\delta(z,x)=\delta(z,m_z)+\delta(m_z,x)$ and similarly for $y$. Let $m=m_a$ and $m'=m_s$. Suppose $m'$ lies on the part of $P$ from $m$ to $y$. Then the path from $s$ to $x$ passes through $m$,
so $\delta(s,x)=\delta(s,m)+\delta(m,x)$, while $\delta(s,a)\le\delta(s,m)+\delta(m,a)$. As $a$ is farthest from $s$, $\delta(s,a)\ge\delta(s,x)$, hence $\delta(m,a)\ge\delta(m,x)$, and
$\delta(a,y)=\delta(a,m)+\delta(m,y)\ge\delta(m,x)+\delta(m,y)=L$. If $m'$ lies on the part from $m$ to $x$, the same argument gives $\delta(a,x)\ge L$. Either way $a$ is the end of a longest path, and the second search finds $\delta(a,b)=L$.
\end{solution}

\begin{solution}{ex:euler}
\idea Even degrees mean that every time you walk into a vertex other than the start, an unused exit is waiting.
\writeup Walking from $v_0$, each visit to a vertex $v\ne v_0$ uses one edge in and one out, so when we arrive at $v$ we have used an odd number of its edges; $\deg v$ being even, one is left. So the walk can only get stuck at $v_0$, and
it is a closed walk. If some edge is unused, connectivity gives an unused edge at a vertex $u$ of the walk (follow a path from $v_0$ to that edge); in the graph of unused edges degrees are still even, so a walk from $u$ closes at $u$ and
can be spliced in. In linear time: keep a stack, initially $(v_0)$; while it is non-empty, look at its top $v$; if $v$ has an unused edge $vw$ (keep a pointer into each adjacency list and mark edges used), mark it and push $w$; otherwise pop $v$
and append it to the output. The output, read backwards, is an Euler circuit: the stack is always a walk from $v_0$, and popping happens exactly when the current sub-walk is closed and exhausted.
\end{solution}

\begin{solution}{ex:countpaths}
\idea A shortest path to $v$ is a shortest path to its last-but-one vertex followed by one edge, and BFS treats vertices level by level.
\writeup Set $N(s)=1$. When $u$ is removed from the queue and scans a neighbour $v$: if $v$ is new, set $d(v)=d(u)+1$, $N(v)=N(u)$; else if $d(v)=d(u)+1$, add $N(u)$ to $N(v)$. By Lemma~\ref{lem:queue}, all $u$ with $d(u)=k$ are removed before any $v$ with
$d(v)=k+1$ is, so $N(u)$ is final when it is used. Shortest paths to $v$ correspond bijectively to pairs (a neighbour $u$ with $d(u)=d(v)-1$, a shortest path to $u$), so $N(v)=\sum N(u)$ over those $u$.
\end{solution}

\begin{solution}{ex:dijkstraneg}
\idea Dijkstra trusts a vertex as soon as it is the closest unsettled one, but a negative edge can later offer a shortcut.
\writeup With $w(s,a)=2$, $w(s,b)=3$, $w(b,a)=-2$: $a$ is extracted and settled with $D(a)=2$, before $b$, but $\delta(s,a)=1$ through $b$. Adding $3$ to every weight gives $5,6,1$: now $s\to a$ costs $5$ and $s\to b\to a$ costs $7$, so the shortest path changed; adding a
constant $c$ adds $kc$ to a path of $k$ edges and favours paths with few edges. (Johnson's reweighting adds $h(u)-h(v)$ instead, which telescopes.)
\end{solution}

\begin{solution}{ex:arbitrage}
\idea Products become sums under logarithms, and ``greater than $1$'' becomes ``negative''.
\writeup Put $w(i,j)=-\log r_{ij}$. A cycle multiplies the money by $\prod r_{i_ti_{t+1}}>1$ if and only if its weight $-\sum\log r$ is negative. Run the Floyd--Warshall recurrence in $O(n^3)$; by Exercise~\ref{ex:fwneg}, there is
an arbitrage if and only if some $D_n(i,i)<0$. (Exhibiting the cycle takes a little more bookkeeping; the notebook does it. Real exchanges also charge fees, which add positive weights.)
\end{solution}

\begin{solution}{ex:dagsp}
\idea In a topological order every path runs forward, so the order of relaxations required by Lemma~\ref{lem:relax}(ii) happens automatically.
\writeup Sort topologically in linear time (Theorem~\ref{thm:cycle}). For every $v$ reachable from $s$, a shortest path exists (there is no cycle at all) and its edges leave vertices in topological order, so they are relaxed in path order:
Lemma~\ref{lem:relax}(ii) gives $D(v)=\delta(s,v)$. Each edge is relaxed once: linear time. For longest paths, negate the weights (still no cycle, so no negative cycle). For a project, add a start and an end vertex, put the duration of each task on
its outgoing edges, and the longest path from start to end is the least possible duration of the project; its edges are the critical tasks.
\end{solution}

\begin{solution}{ex:widest}
\idea Width obeys the same monotonicity as distance: extending a path can only make it narrower, just as it can only make it longer with nonnegative weights.
\writeup Keep $W(s)=+\infty$, $W(v)=0$ otherwise, relax by
\[
W(v)\gets\max\big(W(v),\min(W(u),c(u,v))\big),
\]
and extract the vertex of largest $W$. The proof of Theorem~\ref{thm:dijkstra} goes through with $\le$ reversed: if $u$ is the first vertex
extracted with $W(u)<\omega(u)$, the optimal width, take an optimal path to $u$ and its first edge $(x,y)$ leaving the settled set; then
\[
W(y)\ge\min\big(\omega(x),c(x,y)\big)\ge\omega(u)>W(u),
\]
the middle inequality because the prefix of the optimal path to $u$ is at least as wide as the whole path, while $u$ was extracted with the largest key and $y$ was a candidate: a contradiction. For the second part, let $T$ be a maximum spanning tree and $P$ the path of $T$ from $s$ to $t$, with narrowest edge $e$. If some path from $s$ to $t$ had all capacities $>c(e)$, one of its edges $f$ would reconnect the two parts of $T-e$, and $T-e+f$ would be a heavier spanning tree.
\end{solution}

\begin{solution}{ex:robbins}
\idea Depth-first search already orients the graph: down along the tree, up along back edges. A bridge is the only thing that could stop you from climbing back.
\writeup A bridge, oriented either way, cannot be crossed back, so its two sides cannot reach each other. Conversely, assume no bridge, and orient tree edges from parent to child and back edges from descendant to ancestor (Lemma~\ref{lem:noncross}).
The root reaches everyone by tree edges. Let $v$ be another vertex with parent $u$. Since $uv$ is not a bridge, $\mathrm{low}(v)\le d(u)$ by Theorem~\ref{thm:cutvertices}(3): some back edge goes from a vertex $x$ of $T_v$ to a vertex $w$ with $d(w)\le d(u)$,
a proper ancestor of $v$. Walking down the tree from $v$ to $x$ and then along $x\to w$, $v$ reaches a proper ancestor. Repeating, $v$ reaches the root. Everything is computed by one search.
\end{solution}

\begin{solution}{ex:fwneg}
\idea Entries are weights of walks, so they can only become negative on the diagonal by a negative closed walk; and a negative cycle is caught when its largest vertex is allowed as an intermediate.
\writeup By induction on the recurrence, every finite $D_k(i,j)$ is the weight of a walk from $i$ to $j$; a closed walk splits into cycles, so if no cycle is negative, $D_n(i,i)\ge0$. Conversely, the claim of the hint holds without any assumption, by induction on $k$: if
$k$ is not an intermediate vertex of $p$, $D_k(i,j)\le D_{k-1}(i,j)\le w(p)$; if it is, it cuts $p$ into two paths with distinct intermediates in $\{1,\dots,k-1\}$, and $D_k(i,j)\le D_{k-1}(i,k)+D_{k-1}(k,j)\le w(p)$. Let $C$ be a negative cycle without repeated vertices, $k$
its largest vertex. Read it as a path from $k$ to $k$ whose intermediate vertices are smaller than $k$: $D_{k-1}(k,k)\le w(C)<0$, and entries never increase, so $D_n(k,k)<0$.
\end{solution}

\begin{solution}{ex:zeroone}
\idea A deque holding two consecutive values is a priority queue with two buckets.
\writeup Start with $s$ in the deque. Remove the front $u$ (skip it if already settled, settle it otherwise) and relax its edges, putting the improved $v$ at the front for a $0$-edge, at the back for a $1$-edge. Exactly as in Lemma~\ref{lem:queue}, the keys in the deque are nondecreasing
from front to back and span at most $1$: the front $u$ is minimal, a $0$-edge gives $D(v)=D(u)$, at most the old front, and a $1$-edge gives $D(u)+1$, at least the old back. So each settled vertex is one of least key, and the proof of Theorem~\ref{thm:dijkstra} applies.
Each edge causes at most one insertion, so there are $O(n+m)$ deque operations, each $O(1)$.
\end{solution}

\section*{Sources and further reading}

\begin{itemize}
\item \emph{The stories.} \'E.~Lucas, \emph{R\'ecr\'eations math\'ematiques}, vol.~1, Gauthier-Villars, Paris, 1882, the chapter on mazes, reporting Tr\'emaux's method; S.~Even, \emph{Graph Algorithms}, 2nd ed., Cambridge University Press, 2012, Chapter~3, on
Tr\'emaux and depth-first search. P.~L.~Frana and T.~J.~Misa, \emph{An interview with Edsger W.~Dijkstra}, Communications of the ACM 53(8) (2010), 41--47, for the caf\'e terrace and the ``twenty-minute invention'', paraphrased above;
E.~W.~Dijkstra, \emph{A note on two problems in connexion with graphs}, Numerische Mathematik 1 (1959), 269--271.
\item R.~E.~Tarjan, \emph{Depth-first search and linear graph algorithms}, SIAM Journal on Computing 1 (1972), 146--160; J.~E.~Hopcroft and R.~E.~Tarjan, \emph{Algorithm 447: efficient algorithms for graph manipulation}, Communications of the ACM 16 (1973), 372--378;
M.~Sharir, \emph{A strong-connectivity algorithm and its applications in data flow analysis}, Computers and Mathematics with Applications 7 (1981), 67--72 (Kosaraju's algorithm, from his unpublished 1978 lecture notes, is attributed to him in A.~V.~Aho, J.~E.~Hopcroft and J.~D.~Ullman, \emph{Data Structures and Algorithms}, Addison-Wesley, 1983);
B.~Aspvall, M.~F.~Plass and R.~E.~Tarjan, \emph{A linear-time algorithm for testing the truth of certain quantified boolean formulas}, Information Processing Letters 8 (1979), 121--123.
\item R.~Bellman, \emph{On a routing problem}, Quarterly of Applied Mathematics 16 (1958), 87--90; L.~R.~Ford, \emph{Network flow theory}, RAND paper P-923, 1956; E.~F.~Moore, \emph{The shortest path through a maze}, Proceedings of an International Symposium on the Theory of Switching, Harvard, 1959;
R.~W.~Floyd, \emph{Algorithm 97: Shortest path}, Communications of the ACM 5 (1962), 345; S.~Warshall, \emph{A theorem on Boolean matrices}, Journal of the ACM 9 (1962), 11--12; D.~B.~Johnson, \emph{Efficient algorithms for shortest paths in sparse networks}, Journal of the ACM 24 (1977), 1--13;
M.~L.~Fredman and R.~E.~Tarjan, \emph{Fibonacci heaps and their uses in improved network optimization algorithms}, Journal of the ACM 34 (1987), 596--615.
\item D.~R.~Karger, P.~N.~Klein and R.~E.~Tarjan, \emph{A randomized linear-time algorithm to find minimum spanning trees}, Journal of the ACM 42 (1995), 321--328; B.~Chazelle, \emph{A minimum spanning tree algorithm with inverse-Ackermann type complexity}, Journal of the ACM 47 (2000), 1028--1047;
S.~Pettie and V.~Ramachandran, \emph{An optimal minimum spanning tree algorithm}, Journal of the ACM 49 (2002), 16--34; H.~E.~Robbins, \emph{A theorem on graphs, with an application to a problem of traffic control}, American Mathematical Monthly 46 (1939), 281--283;
C.~Hierholzer, \emph{\"Uber die M\"oglichkeit, einen Linienzug ohne Wiederholung und ohne Unterbrechung zu umfahren}, Mathematische Annalen 6 (1873), 30--32.
\item T.~H.~Cormen, C.~E.~Leiserson, R.~L.~Rivest and C.~Stein, \emph{Introduction to Algorithms}, 4th ed., MIT Press, 2022, Chapters~20--23, follow the same route with the same proofs of the parenthesis and white-path theorems.
\end{itemize}
